\section{Statistical Inference: Estimation}

\subsection{Introduction to Statistical Inference}
Statistical inference is the art of drawing conclusions about a whole population based on a smaller sample. Since we operate under \textbf{uncertainty}, we use probability theory to quantify how much we can trust our conclusions.

\begin{defbox}{Core Definitions}
\begin{itemize}
    \item \textbf{Statistical Inference:} Drawing conclusions about a population from a sample.
    \item \textbf{Sampling Observation:} A single observed value of a random variable.
\end{itemize}
\end{defbox}


\subsubsection{Branches of Inference}
There are two main ways to perform inference:
\begin{enumerate}
    \item \textbf{Estimation Theory:} Determining the value of unknown parameters.
    \begin{itemize}
        \item \textit{Point Estimation:} A single \textit{best guess} value.
        \item \textit{Confidence Interval:} A range of likely values with a associated probability.
    \end{itemize}
    \item \textbf{Testing Theory (Hypothesis Testing):} Empirically verifying a claim about a parameter or a model (e.g., "Is this new feature increasing user clicks?").
\end{enumerate}

\subsubsection{Frequentist vs. Bayesian Statistics}
This is the most fundamental philosophical divide in Data Science. The difference lies in how they treat the unknown parameter $\theta$.

\begin{itemize}
    \item \textbf{Classical (Frequentist) Statistics:} $\theta$ is a \textbf{fixed but unknown constant}. Uncertainty comes from \textit{sampling variability} (if I took a different sample, I'd get a different result). We only use the current data.
    \item \textbf{Bayesian Statistics:} $\theta$ is a \textbf{random variable}. We start with \textbf{Prior} beliefs, observe data (\textbf{Likelihood}), and update our knowledge into a \textbf{Posterior} distribution.
\end{itemize}

\begin{defbox}{Bayes' Theorem for Inference}
\[
\underbrace{P(\theta | \text{data})}_{\text{Posterior}} = \frac{\overbrace{P(\text{data} | \theta)}^{\text{Likelihood}} \times \overbrace{P(\theta)}^{\text{Prior}}}{P(\text{data})}
\]
\end{defbox}


\vs

\begin{example} {\textbf{An Analogy:}}
  
  Imagine guessing a person's age without seeing them.

\textit{Frequentist:} "This person has a true age.  Say, 35. I don't know it, but it's an objective fact. I can only estimate it from clues."

\textit{Bayesian:} "Before any information, I think they're between 20–60 years old (vague prior). Then I learn they've worked for 15 years $\rightarrow$ I update to 35–50. Then I see they have gray hair $\rightarrow$ I update again to 45–55."

The Bayesian continuously \textit{updates beliefs} as new evidence arrives, while the Frequentist focuses on \textit{what the data alone tells us} about a fixed truth.

\end{example}

\subsection{Fundamental Concepts of Estimation}

\subsubsection{Estimator vs. Estimate}
In Data Science, we often confuse the \textit{algorithm} with the \textit{result}.
\begin{itemize}
    \item \textbf{Estimator ($T_n$):} It is the \textbf{formula} (a random variable). For example, the function `mean()`. It depends on the sample, so it changes every time you pick a new dataset.
    \item \textbf{Estimate ($t_n$):} It is the \textbf{number} (a fixed value). For example, "173 cm". It is the result of applying the formula to your specific data.
\end{itemize}

\subsubsection{The Sampling Distribution}
The \textbf{Sampling Distribution} is a theoretical histogram of all possible estimates you would get if you could repeat your experiment infinite times. It is the key to understanding the \textbf{uncertainty} of your model.

\subsubsection{Measuring Quality: Mean Squared Error (MSE)}


To evaluate the quality of an estimator, we need a measure of how \textit{close} the estimator is to the true parameter value on average.

\begin{defbox}{Estimation Error}
The Estimation error of an estimator $T_n$ for parameter $\theta$ is:
\begin{equation}
\text{err}(T_n) = T_n - \theta
\end{equation}
\end{defbox}

\begin{defbox}{Mean Squared Error}
The Mean Squared Error (MSE) of an estimator $T_n$ is defined as:
\begin{equation}
\text{MSE}(T_n) = E\left[(T_n - \theta)^2\right]
\end{equation}
This measures the expected squared distance between the estimator and the true parameter value.
\end{defbox}

\textbf{Why use MSE?} 
1. It ignores the sign (errors of +5 and -5 are both bad).
2. It penalizes \textbf{outliers} (large errors) more heavily because they are squared.

\subsubsection{The Golden Rule: Bias-Variance Decomposition}
This is the most important formula for your career.  When evaluating if an algorithm is good, we look at its \textbf{Mean Squared Error (MSE)}. The magical property of MSE is that it mathematically breaks down into two distinct sources of error:


\begin{defbox}{MSE Decomposition}
\begin{equation}
\text{MSE}(T_n) = \text{Var}(T_n) + \left[\text{bias}(T_n)\right]^2
\end{equation}
\end{defbox}




\begin{itemize}
    \item \textbf{Bias (Systematic Error):  $\text{bias}(T_n) = E[T_n] - \theta$  :} Does the model always miss the target in the same direction? (e.g., a scale that always adds 1kg).
    \item \textbf{Variance (Sensitivity to data): $\text{Var}(T_n) = E\left[(T_n - E[T_n])^2\right]$} Is the model unstable? Does it change drastically if you change a few rows in your dataset?
\end{itemize}

\pagebreak
\textbf{Data Science Insight:} 
\begin{itemize}
    \item \textbf{High Bias = Underfitting:} The model is too simple (it misses the pattern).
    \item \textbf{High Variance = Overfitting:} The model is too complex (it memorizes the noise in the data).
\end{itemize}

\vs

{\footnotesize
\textbf{Proof of the MSE Decomposition:}

Starting from the definition of MSE:
\begin{align*}
\text{MSE}(T_n) &= E\left[(T_n - \theta)^2\right]
\end{align*}

We add and subtract $E[T_n]$ inside the square:
\begin{align*}
&= E\left[\left((T_n - E[T_n]) + (E[T_n] - \theta)\right)^2\right]
\end{align*}

Expanding the square $(a + b)^2 = a^2 + 2ab + b^2$:
\begin{align*}
&= E\left[(T_n - E[T_n])^2\right] + E\left[(E[T_n] - \theta)^2\right] + 2E\left[(T_n - E[T_n])(E[T_n] - \theta)\right]
\end{align*}

Now we analyze each term:
\begin{itemize}
    \item First term: $E\left[(T_n - E[T_n])^2\right] = \text{Var}(T_n)$
    \item Second term: $(E[T_n] - \theta)^2 = [\text{bias}(T_n)]^2$ (this is a constant, so $E[\text{constant}] = \text{constant}$)
    \item Third term: Since $(E[T_n] - \theta)$ is a constant, we can factor it out:
    \[
    2(E[T_n] - \theta) \cdot E[T_n - E[T_n]] = 2(E[T_n] - \theta) \cdot 0 = 0
    \]
\end{itemize}

Therefore:
\[
\text{MSE}(T_n) = \text{Var}(T_n) + [\text{bias}(T_n)]^2
\]
}

\vspace{0.2cm}
\begin{example}{\textbf{Computing MSE with Bias-Variance Decomposition}}

Suppose we estimate the population mean $\mu = 100$ using an estimator $T_n$ that has:
\begin{itemize}
    \item Expected value $E[T_n] = 102$ (so it overestimates on average).
    \item Variance $\text{Var}(T_n) = 9$.
\end{itemize}

\begin{itemize}
    \item \texttt{Step 1:} Calculate the bias:
    \[
    \text{bias}(T_n) = E[T_n] - \theta = 102 - 100 = 2
    \]
    \item \texttt{Step 2:} Calculate the squared bias:
    \[
    [\text{bias}(T_n)]^2 = 2^2 = 4
    \]
    \item \texttt{Step 3:} Calculate the MSE using the decomposition:
    \[
    \text{MSE}(T_n) = \text{Var}(T_n) + [\text{bias}(T_n)]^2 = 9 + 4 = 13
    \]
\end{itemize}

This decomposition shows that part of the error comes from the estimator's variability (variance = 9) and part from its systematic tendency to overestimate (squared bias = 4).
\end{example}
    

\subsubsection{BLUE (Best Linear Unbiased Estimator)}
In many algorithms, particularly Ordinary Least Squares (OLS) regression, we look for an estimator that is \textbf{BLUE}. To earn this title, the estimator must strictly meet three conditions:

\begin{defbox}{The 3 Conditions of BLUE}
\begin{enumerate}
    \item \textbf{Linear:} It is calculated as a simple linear combination of the sample data.
    \item \textbf{Unbiased:} On average, it perfectly hits the true parameter ($E[T_n] = \theta$).
    \item \textbf{Best (Most Efficient):} It has the \textit{lowest possible variance} among all other estimators that are both linear and unbiased.
\end{enumerate}
\end{defbox}

\begin{defbox}{Unbiased Estimator}
An estimator $T_n$ is unbiased if and only if:
\[
E[T_n] = \theta
\]
In other words, on average (across all possible samples), the estimator equals the true parameter value.
\end{defbox}

An unbiased estimator has no systematic tendency to overestimate or underestimate the parameter. This is often considered a desirable property, though as we'll see, it's not always the most important criterion.


    \begin{figure}[H] 
        \centering\includegraphics[scale=0.35]{139.Unbiased Estimator.jpeg}
        \caption{Unbiased Estimator}\captionsetup{font=footnotesize}
        \label{fig - Unbiased Estimator}
    \end{figure}




\textit{Data Science Takeaway:} If an estimator is BLUE, you can be mathematically certain that no other linear equation will give you a more precise, unbiased prediction.


\begin{defbox}{Relative Efficiency}
An estimator $T_n$ is relatively more efficient than another estimator $T_n'$ if and only if:
\[
\text{MSE}(T_n) < \text{MSE}(T_n')
\]
For unbiased estimators, this simplifies to comparing variances: $\text{Var}(T_n) < \text{Var}(T_n')$.
\end{defbox}

\begin{defbox}{Most Efficient Estimator (Absolute Efficiency)}
An estimator $T_n$ is the most efficient if and only if:
\[
\text{MSE}(T_n) = \kappa(\theta)
\]
where $\kappa(\theta)$ is the minimum achievable MSE for any estimator of $\theta$ under specific regularity conditions. This lower bound is known as the \textbf{Fréchet-Rao-Cramér bound} (1943-1946).
\end{defbox}

    \begin{figure}[H] 
        \centering\includegraphics[scale=0.4]{140.Efficiency.jpeg}
        \caption{Efficiency}\captionsetup{font=footnotesize}
        \label{fig - Efficiency}
    \end{figure}


\begin{defbox}{Sufficient Estimator}
An estimator $T_n$ is sufficient when it gives as much information about $\theta$ as is possible from the sample at hand. Once we know the value of a sufficient statistic, no other function of the sample can provide additional information about $\theta$.
\end{defbox}


\subsubsection{UMVUE (Uniformly Minimum Variance Unbiased Estimator)}
While BLUE restricts the competition only to \textit{linear} estimators, \textbf{UMVUE} is the absolute gold standard across the board. 

\begin{defbox}{UMVUE}
An estimator is UMVUE if it has the smallest possible variance among \textbf{ALL} unbiased estimators (whether they are linear or non-linear). The term \textit{Uniformly} means this superiority holds true for every single possible value of the parameter $\theta$.
\end{defbox}

If you find a UMVUE for your model, you have reached the theoretical limit of precision for an unbiased estimator. It is mathematically impossible to find another unbiased estimator that is consistently more accurate.

\begin{figure}[H] 
    \centering\includegraphics[scale=0.4]{141.UMVE.jpeg}
    \caption{UMVUE: Among all the unbiased estimators (all perfectly centered on $\theta$), the UMVUE has the narrowest and most concentrated distribution.}
    \captionsetup{font=footnotesize}
    \label{fig - UMVE}
\end{figure}

\subsection{Common Estimators}

This section presents the three most frequent estimators you will use to summarize datasets. Understanding their expected values and variances is the first step in hypothesis testing.

\subsubsection{Sample Mean}
The sample mean is the standard estimator for the population mean $\mu$. For most theoretical probability models (Normal, Uniform, Exponential), it is the \textbf{BLUE} (Best Linear Unbiased Estimator).

\begin{defbox}{Sample Mean}
\begin{equation}
\bar{X}_n = \frac{1}{n}\sum_{i=1}^{n} X_i
\end{equation}
\textbf{Key Properties:}
\begin{itemize}
    \item \textbf{Expected Value:} $E[\bar{X}_n] = \mu$ (It is unbiased).
    \item \textbf{Variance:} $\text{Var}(\bar{X}_n) = \frac{\sigma^2}{n}$ (The precision increases as the sample size $n$ grows).
\end{itemize}
\end{defbox}

\subsubsection{Sample Proportion}
When dealing with binary events (e.g., clicks vs. no-clicks, churn vs. retention), your data follows a Bernoulli distribution. The estimator for the probability of success $\pi$ is the sample proportion.

\begin{defbox}{Sample Proportion}
\begin{equation}
P_n = \frac{1}{n}\sum_{i=1}^{n} X_i \quad \text{where } X_i \in \{0, 1\}
\end{equation}
\textbf{Key Properties:}
\begin{itemize}
    \item \textbf{Expected Value:} $E[P_n] = \pi$ (It is unbiased).
    \item \textbf{Variance:} $\text{Var}(P_n) = \frac{\pi(1-\pi)}{n}$
\end{itemize}
\end{defbox}

\textit{Data Science Takeaway:} If you survey $100$ users and $65$ click a button, your estimate is $P_{100} = 0.65$. The Standard Error of this estimate is simply the square root of its variance: $\text{SE} = \sqrt{\frac{0.65 \times 0.35}{100}} \approx 0.047$. 

\subsubsection{Sample Variance ($n$ vs $n-1$)}
Estimating how much your data varies is tricky. There are two different formulas, and choosing the wrong one can lead to underestimating the risk or spread in your data.

\begin{defbox}{Biased Sample Variance ($S_n^2$)}
\begin{equation}
S_n^2 = \frac{1}{n}\sum_{i=1}^{n} (X_i - \bar{X}_n)^2
\end{equation}
This formula is \textbf{biased}. It systematically underestimates the true population variance $\sigma^2$ because $E[S_n^2] = \frac{n-1}{n}\sigma^2$.
\end{defbox}

\begin{defbox}{Corrected Unbiased Sample Variance ($\tilde{S}_n^2$)}
To fix the underestimation, we apply \textbf{Bessel's Correction} by dividing by $n-1$ instead of $n$:
\begin{equation}
\tilde{S}_n^2 = \frac{1}{n-1}\sum_{i=1}^{n} (X_i - \bar{X}_n)^2
\end{equation}
This version is \textbf{unbiased}, meaning $E[\tilde{S}_n^2] = \sigma^2$. 
\end{defbox}

\vs
\textbf{Why is dividing by $n$ biased?} \\
When you calculate variance, you measure how far each data point is from the center. Ideally, you would measure the distance from the \textit{true population mean} $\mu$. However, because $\mu$ is unknown, you are forced to measure the distance from the \textit{sample mean} $\bar{X}_n$. 

Because $\bar{X}_n$ is calculated directly from your sample, it sits exactly in the mathematical center of your specific data points. This makes the measured distances artificially small. Dividing by $n-1$ perfectly corrects for this artificial mathematical shrinking.

\textit{Note:} The relationship between the two is simply $\tilde{S}_n^2 = \frac{n}{n-1} S_n^2$. As your dataset grows to millions of rows, the difference between $n$ and $n-1$ becomes mathematically irrelevant.

\subsection{Asymptotic Properties of Estimators (Large Samples)}

When dealing with massive datasets ($n \to \infty$), we care less about how an estimator behaves on a tiny sample and more about its \textbf{asymptotic} behavior. Does it eventually converge to the truth?

\vspace{-0.2cm}
\paragraph*{1. Asymptotic Unbiasedness.}
An estimator is \textbf{asymptotically unbiased} if any systematic error (bias) disappears as your dataset grows to infinity.

\begin{defbox}{Asymptotic Unbiasedness}
\begin{equation}
\lim_{n \to \infty} E[T_n] = \theta
\end{equation}
\end{defbox}

\textit{Data Science Takeaway:} This is why we don't panic over the biased sample variance ($S_n^2$). While it underestimates the variance slightly in small samples, the difference between dividing by $n$ and $n-1$ becomes mathematically zero when $n = 1,000,000$. 


\vspace{-0.2cm}
\paragraph*{2. Asymptotic Efficiency.}
An estimator is asymptotically efficient if, given infinite data, its Mean Squared Error (MSE) drops to exactly zero. 

\begin{defbox}{Asymptotic Efficiency}
\[ \lim_{n \to \infty} \text{MSE}(T_n) = 0 \]
Since MSE is the sum of Variance and squared Bias, this means both the bias \textit{and} the variance must disappear entirely as $n \to \infty$.
\end{defbox}


\vspace{-0.2cm}
\paragraph*{3. Consistency.}
Consistency is the single most important property in all of Big Data analytics. An estimator is consistent if gathering more data guarantees that your estimate will eventually lock onto the true parameter value.

\begin{defbox}{Consistency}
An estimator $T_n$ is \textbf{weakly consistent} if the probability of missing the true value ($\theta$) by even a microscopic amount ($\varepsilon$) shrinks to zero:
\begin{equation}
\lim_{n \to \infty} P(|T_n - \theta| < \varepsilon) = 1
\end{equation}
\end{defbox}

\begin{figure}[H] 
    \centering\includegraphics[scale=0.5]{142.Consistency.jpeg}
    \caption{Consistency: As $n$ grows (20, 50, 100, 500), the distribution collapses from a wide, uncertain bell curve into a perfectly sharp spike directly on the true value $\theta$.}
    \captionsetup{font=footnotesize}
    \label{fig - Consistency}
\end{figure}

\subsubsection{The Law of Large Numbers (LLN)}
The Law of Large Numbers is the mathematical engine that drives consistency. It proves that the sample mean ($\bar{X}_n$) is a consistent estimator of the population mean ($\mu$). 

\begin{defbox}{The Law of Large Numbers}
\begin{itemize}
    \item \textbf{Weak LLN (Convergence in Probability):} $\bar{X}_n \xrightarrow{P} \mu$. The probability of your sample average being wrong goes to zero.
    \item \textbf{Strong LLN (Convergence Almost Surely):} The sequence of sample averages will, with 100\% certainty, hit the exact expected value $\mu$ given infinite trials.
\end{itemize}
\end{defbox}

\vspace{0.2cm}
\begin{example}{\textbf{The Casino Advantage}}

If you roll a fair die, the mathematical expected value is $3.5$. 
\begin{itemize}
    \item If you roll it 10 times, your average might be $2.9$ or $4.1$. (High variance).
    \item If you roll it $10,000$ times, your average will be exactly $3.501$. 
\end{itemize}
This is the Law of Large Numbers in action. It is the mathematical reason why casinos always make money in the long run, regardless of what happens on any individual roulette spin.
\end{example}

\subsection{Limit Theorems}

Limit theorems describe what happens to data when your sample size $n$ grows to infinity. They are the mathematical justification for why we can use the Normal distribution to solve real-world problems, even when the underlying data is a completely chaotic mess.

\subsubsection{The Central Limit Theorem (CLT)}

This is arguably the single most important theorem in all of Data Science.

\begin{defbox}{Central Limit Theorem}
Let $X_1, X_2, \ldots, X_n$ be a sample of $n$ independent and identically distributed (i.i.d.) random variables with mean $\mu$ and finite variance $\sigma^2$. 

\vs

As $n \to \infty$, the distribution of the \textbf{sample mean} converges to a Normal distribution:
\begin{equation} \bar{X}_n \xrightarrow{d} N\left(\mu, \frac{\sigma^2}{n}\right) \end{equation}
In its standardized form (Z-score):
\begin{equation}\frac{\bar{X}_n - \mu}{\sigma/\sqrt{n}} \xrightarrow{d} N(0, 1) \end{equation}
\end{defbox}

\textbf{Why is this a superpower?} The theorem holds \textbf{regardless of the original shape of the data}. Your users' session times might follow an Exponential distribution, and your sales might follow a Gamma distribution. It doesn't matter. If you take a large enough sample ($n \geq 30$ is usually the rule of thumb), the distribution of their \textit{averages} will always be a perfect Gaussian bell curve.

  



\subsubsection{De Moivre-Laplace Theorem (CLT for Proportions)}

What if your data isn't continuous numbers, but simple Yes/No binary outcomes? (e.g., Did the user click the ad? Yes=1, No=0). The De Moivre-Laplace theorem is just the Central Limit Theorem specifically applied to \textbf{Bernoulli} variables.

\begin{defbox}{De Moivre-Laplace Theorem}
Let $X_1, X_2, \ldots, X_n$ be $n$ Bernoulli random variables with success probability $\pi$ and variance $\pi(1-\pi)$. 

\vs
As $n \to \infty$, the \textbf{sample proportion} (the average of the 1s and 0s) converges to a Normal distribution:
$$ \bar{X}_n \xrightarrow{d} N\left(\pi, \frac{\pi(1-\pi)}{n}\right) $$

In standardized form:
$$ \frac{\bar{X}_n - \pi}{\sqrt{\pi(1-\pi)/n}} \xrightarrow{d} N(0, 1) $$
\end{defbox}


  \begin{figure}[H] 
        \centering\includegraphics[scale=0.5]{143.Central Limit Theorem.jpeg}
        \caption{Central Limit Theorem }
        \captionsetup{font=footnotesize}
        \label{fig - Central Limit Theorem}
    \end{figure}

    Notice how the chaotic shapes at the top all morph into perfect normal distributions as $n$ increases.
   
\textit{Data Science Takeaway:} This theorem is the mathematical engine behind \textbf{A/B Testing}. When you are comparing the Conversion Rates of two different landing pages, you are relying on the De Moivre-Laplace theorem to assume your conversion rates are normally distributed, allowing you to calculate p-values.


    \subsection{Estimation Methods}

This section introduces the main approaches for deriving estimators. Each method provides a systematic way to construct estimators with desirable properties. The main Estimations methods are: 
\begin{itemize}
    \item Least Squares Estimation (LSE).
    \item Maximum Likelihood Estimation (MLE). 
    \item Bayesian Estimation. 
\end{itemize}

\subsubsection{Least Squares Estimation (LSE)}
Used heavily in Linear Regression. It completely ignores probability distributions. It simply draws a line that minimizes the sum of the squared physical distances between the predictions and the actual data points.
\begin{equation}
    \min_{\theta} \sum_{i=1}^{n} \left(x_i - g(\theta)\right)^2
\end{equation}


\vspace{0.2cm}
\begin{example}{\textbf{Understanding Least Squares Estimation:}}


Imagine you have a scatter plot. For example, the horizontal axis could be the \textit{square footage of a house} and the vertical axis the \textit{sale price}. The points on the graph represent actual houses that have been sold.

Your goal in \textbf{Linear Regression} is to draw a straight line that passes through these points in the best possible way, allowing you to guess the price of future houses. But how do you decide which is the \textbf{best} line among the infinite lines you could potentially draw?

This is where \textbf{Least Squares Estimation (LSE)} comes into play. The text highlights two important concepts:
\begin{itemize}
    \item \textbf{It completely ignores probability distributions}: There are other statistical methods (like \textit{Maximum Likelihood Estimation} or MLE) that find the best line based on complex assumptions, such as \textit{assuming the errors follow a Gaussian bell curve}. LSE totally ignores this. It is a purely \textbf{geometric and physical} approach. It literally measures the physical distance (as if using a ruler) between the actual data point on the graph and the line you drew.
    \item \textbf{Minimizes the sum of the squared physical distances}: The best line, according to this method, is simply the one that makes the smallest total error.
\end{itemize}



\end{example}



\subsubsection{Maximum Likelihood Estimation (MLE)}
Used heavily in Logistic Regression and Neural Networks. It asks a reverse-engineered question: \textit{"Given the data I just observed, what parameter value $\theta$ makes this data the most probable to have occurred?"}

Instead of distances, it maximizes the \textbf{Likelihood Function} $\ell(\theta)$. Because multiplying thousands of probabilities causes computers to underflow (numbers get too small), Data Scientists always maximize the \textbf{Log-Likelihood} instead:
\begin{equation}
    \hat{\theta}_{ML} = \arg\max_{\theta} \sum_{i=1}^{n} \log f(x_i; \theta)
\end{equation}
\textit{Fun Fact:} Minimizing \textit{Cross-Entropy Loss} in a Neural Network is mathematically identical to performing Maximum Likelihood Estimation!

    \begin{figure}[H] 
        \centering\includegraphics[scale=0.5]{145.MLE for Bernoulli Parameter.jpeg}
        \caption{MLE for Bernoulli Parameter}\captionsetup{font=footnotesize}
        \label{fig - MLE for Bernoulli Parameter}
    \end{figure}


\vspace{0.2cm}
\begin{example}{\textbf{Understanding Maximum Likelihood Estimation:}}

\begin{itemize}
    \item \textbf{The Core Idea:} Imagine you have a coin and want to figure out if it's biased. You flip it 10 times and get: \textbf{7 heads and 3 tails}.
\begin{itemize}
    \item The question is: what is the probability $\theta$ that this coin lands heads?
\end{itemize}

\vspace{0.3cm}
\item \textbf{The Likelihood Reasoning:} MLE reasons \textit{backwards}:
\begin{center}
\textit{``I observed this data. Which value of $\theta$ makes it \textbf{most probable} to have observed exactly what I saw?''}
\end{center}
\item Let's try different values of $\theta$:
\begin{itemize}
    \item \texttt{Step 1:} If $\theta = 0.5$ (fair coin):
    \[
    L(0.5) = (0.5)^7 \times (0.5)^3 = (0.5)^{10} = 0.00098
    \]
    \item \texttt{Step 2:} If $\theta = 0.7$ (biased toward heads):
    \[
    L(0.7) = (0.7)^7 \times (0.3)^3 = 0.0822 \times 0.027 = 0.00222
    \]
    \item \texttt{Step 3:} If $\theta = 0.9$ (heavily biased):
    \[
    L(0.9) = (0.9)^7 \times (0.1)^3 = 0.478 \times 0.001 = 0.00048
    \]
    \begin{itemize}
        \item The value $\theta = 0.7$ gives the highest likelihood! Indeed, the MLE is exactly $\hat{\theta} = 7/10 = 0.7$.
    \end{itemize}
\end{itemize}
\vspace{0.3cm}
\item \textbf{Why Use the Log-Likelihood?} The likelihood is a \textbf{product} of many terms:
\[
L(\theta) = f(x_1;\theta) \times f(x_2;\theta) \times \cdots \times f(x_n;\theta)
\]
\item \textit{Practical problems:}
\begin{itemize}
    \item With many data points, the product becomes a tiny number (numerical underflow)
    \item Differentiating a product is complicated
\end{itemize}
\vs

\item  \textit{Solution:} Use the \textbf{logarithm}, which transforms products into sums:
\[
\ell(\theta) = \log L(\theta) = \sum_{i=1}^{n} \log f(x_i; \theta)
\]
Since $\log$ is a monotonically increasing function, maximizing $L(\theta)$ is equivalent to maximizing $\ell(\theta)$.

\vspace{0.3cm}
\item \textbf{Key Difference: Least Squares vs. Maximum Likelihood}

\begin{center}
\begin{tabular}{|l|c|c|}
\hline
\textbf{Aspect} & \textbf{Least Squares} & \textbf{Maximum Likelihood} \\
\hline
Requires distribution? & No & Yes \\
\hline
What does it optimize? & Sum of squared errors & Probability of the data \\
\hline
Assumptions & Only on mean \& variance of errors & Full distribution \\
\hline
\end{tabular}
\end{center}
\end{itemize}

\end{example}


%%%%%%%%%%%%%%%%%%%%%%%%%%%%%%%%%%%%%%%%%%%%%%%%%%%%%%%%%%%%%%%%%%%%%
\subsubsection{Bayesian Estimation}

Bayesian estimation provides a powerful alternative to Maximum Likelihood Estimation (MLE) by incorporating \textbf{prior beliefs} about the parameter before observing any data.

\begin{defbox}{Bayesian Estimation Framework}

\textbf{The Setup:}
\begin{itemize}
    \item Let $X \sim f(x; \theta)$ be the probability model with an unknown parameter $\theta$.
    \item Unlike MLE, we treat $\theta$ as a \textbf{random variable} with its own \textbf{prior probability distribution}, denoted as $g(\theta)$.
    \item Let $C_n = (X_1, \ldots, X_n)$ be an $n$-dimensional independent and identically distributed (i.i.d.) random sample (our observed data).
\end{itemize}

\vspace{0.2cm}
\textbf{The Core Equation (Bayes' Theorem):}
$$P(\theta | \text{data}) = \frac{\ell(\text{data} | \theta) \cdot g(\theta)}{P(\text{data})}$$

\begin{itemize}
    \item \textbf{Posterior} $P(\theta | \text{data})$: What we believe about $\theta$ \textit{after} seeing the data.
    \item \textbf{Likelihood} $\ell(\text{data} | \theta)$: The probability of observing this data given a specific $\theta$.
    \item \textbf{Prior} $g(\theta)$: What we believed about $\theta$ \textit{before} seeing the data.
    \item \textbf{Marginal Likelihood} $P(\text{data})$: A normalizing constant to ensure probabilities sum to 1.
\end{itemize}

\vspace{0.2cm}
\textbf{The Estimator:}
\begin{itemize}
    \item The most common Bayesian point estimate is found by maximizing the posterior probability (known as Maximum A Posteriori, or MAP):
\end{itemize}
$$\max_{\theta \in \Theta} P(\theta | \text{data}) \propto g(\theta) \cdot \ell(\text{data} | \theta)$$

\end{defbox}

\vspace{0.2cm}
\begin{example}{\textbf{Understanding Bayesian Estimation}}

Bayesian estimation is fundamentally about \textbf{updating your beliefs} based on new evidence.
\begin{center}
\textbf{Prior} (Initial Belief) $+$ \textbf{Likelihood} (New Evidence) $\Longrightarrow$ \textbf{Posterior} (Updated Belief)
\end{center}

\vspace{0.3cm}
\textbf{1. The Key Difference: Frequentist vs. Bayesian}
\begin{center}
\renewcommand{\arraystretch}{1.3}
\begin{tabular}{|l|l|}
\hline
\textbf{Frequentist (MLE)} & $\theta$ is a \textbf{fixed} (but unknown) constant. \\
\hline
\textbf{Bayesian} & $\theta$ is a \textbf{random variable} with a probability distribution. \\
\hline
\end{tabular}
\end{center}

\vspace{0.3cm}
\textbf{2. Numerical Example: The Coin Flip}
\begin{itemize}
    \item \textbf{Goal:} Estimate $\pi$ (the true probability of flipping heads).
    \item \textbf{Step 1 (Prior):} You strongly believe the coin is fair. We can represent this mathematically as a distribution centered at 0.5 (e.g., equivalent to having previously seen 10 heads and 10 tails).
    \item \textbf{Step 2 (Data/Likelihood):} You flip the coin 10 times and observe \textbf{8 heads} and \textbf{2 tails}.
    \item \textbf{Step 3 (Estimation):} Let's compare the two approaches:
    \begin{itemize}
        \item \textbf{MLE (Ignores prior):} Looks strictly at the new data.
        $$\hat{\pi}_{ML} = \frac{8 \text{ (heads)}}{10 \text{ (total flips)}} = 0.80$$
        
        \item \textbf{Bayesian (Combines prior + data):} Adds the new data to your initial beliefs.
        $$\hat{\pi}_{Bayes} = \frac{10 \text{ (prior heads)} + 8 \text{ (new heads)}}{20 \text{ (prior total)} + 10 \text{ (new total)}} = \frac{18}{30} = 0.60$$
    \end{itemize}
    \item \textbf{Shrinkage:} The Bayesian estimate ($0.6$) is a compromise. The strong prior "pulled" or "shrank" the final estimate away from the extreme data observation ($0.8$) back toward the initial belief ($0.5$).
\end{itemize}

\vspace{0.3cm}
\textbf{3. When Does the Prior Matter Less?}
\begin{itemize}
    \item The influence of the prior depends entirely on the \textbf{sample size ($n$)}. As you collect more data, the evidence overwhelms your initial beliefs.
\end{itemize}

\begin{center}
\renewcommand{\arraystretch}{1.2}
\begin{tabular}{|c|c|c|l|}
\hline
\textbf{Sample Size ($n$)} & \textbf{Prior Influence} & \textbf{Data Influence} & \textbf{Result} \\
\hline
Small ($n=10$) & High & Low & Prior dominates (Shrinkage is strong) \\
Medium ($n=100$) & Medium & Medium & Balanced compromise \\
Large ($n \to \infty$) & $\approx 0$ & $\approx 100\%$ & \textbf{Bayesian converges to MLE} \\
\hline
\end{tabular}
\end{center}

\end{example}

\pagebreak




\subsection{Confidence Intervals}
A Confidence Interval gives a range of plausible values for a parameter. 

 In Frequentist statistics, the true parameter is a fixed, immutable number. Your specific interval has already been calculated. Therefore, the parameter is either 100\% inside your interval, or 0\% inside it. It doesn't "jump" in and out.
 
 "95\% Confidence" does \textit{not} mean there is a 95\% chance the true parameter is in your specific interval. It means that if you repeated your experiment 100 times, 95 of the intervals you generate would successfully catch the true parameter.

 \begin{defbox}{The Fishing Analogy:} 
\begin{itemize}
    \item A Point Estimate (Sample Mean) is like trying to catch a specific fish with a spear. It's very easy to miss.
    \item A Confidence Interval is like throwing a net. Saying "95\% Confidence" means you are confident in the \textit{reliability of your net-throwing method}, not the probability of the fish moving into the net after it's been thrown.
\end{itemize}
\end{defbox}

\subsubsection{Estimating the Mean: Z-Interval vs. T-Interval}

When building an interval around the sample mean, you must decide how wide to cast your \textit{net} (the margin of error). The choice between $Z$ and $T$ depends entirely on the level of uncertainty in your data.

\subsubsection{The Z-Interval (Known Variance or Large Sample)}
If you know the population variance, or if your sample is large enough ($n > 30$) to invoke the Central Limit Theorem, you use the Standard Normal ($Z$) distribution:
\begin{equation}
    \bar{x}_n \pm z_{\alpha/2}\frac{\sigma}{\sqrt{n}}
\end{equation}

    \begin{figure}[H] 
        \centering\includegraphics[scale=0.38]{146.Confidence Interval sampling mean.jpeg}
        \caption{Confidence Interval sampling mean}\captionsetup{font=footnotesize}
        \label{fig - Confidence Interval sampling mean}
    \end{figure}

\subsubsection{The T-Interval (Unknown Variance \& Small Sample)}
If you do not know the population variance and have a small sample ($n \leq 30$), you must estimate the variance using the sample standard deviation ($s$). Because this adds uncertainty, you must use the \textbf{Student's t-distribution}, which has heavier tails to widen the interval and protect against error.
\begin{equation}
    \bar{x}_n \pm t_{n-1, \alpha/2}\frac{s}{\sqrt{n}}
\end{equation}


\subsubsection{Confidence Intervals for the Variance ($\chi^2$)}

The $Z$ and $T$ intervals are used to estimate the \textit{center} of the data (the mean). But what if you want to estimate how spread out the data is by building an interval for the \textbf{variance itself} ($\sigma^2$)?

You \textbf{cannot} use $Z$ or $T$. The mathematical reason is simple:
\begin{itemize}
    \item The $Z$ and $T$ curves are perfectly \textbf{symmetric} and centered at zero. They include negative numbers.
    \item Variance ($\sigma^2$) is a squared value. \textbf{It cannot physically be negative.} The absolute minimum variance is zero (which happens only if all data points are identical).
\end{itemize}

Because variance "hits a wall" at zero and can only stretch towards positive infinity, its distribution is not a symmetric bell. It is heavily \textbf{right-skewed}. 

To model this, we use the \textbf{Chi-Squared ($\chi^2$)} distribution. Because this curve is asymmetrical, you cannot use the $\pm$ symbol (which adds and subtracts the exact same margin). Instead, you must calculate the left boundary and the right boundary using two completely different values from the Chi-Squared table, resulting in an \textbf{asymmetric confidence interval}.


    \begin{figure}[H] 
        \centering\includegraphics[scale=0.4]{147.Confidence Interval of the popultation variance.jpeg}
        \caption{Confidence Interval of the popultation variance}\captionsetup{font=footnotesize}
        \label{fig - Confidence Interval of the popultation variance}
    \end{figure}
